\chapter*{Analysing Results from Molecular Docking}

In addition to advanced rescoring methods, it is also important to visually inspect the docking solution to assess if the protein-ligands show prominent interactions such as hydrogen bonds, salt bridges, water-mediated interactions, cation-$\pi$ and $\pi$-$\pi$. Since the results need to be communicated to a wider audience, the colour coding for receptor and ligand atoms must be distinct. The atom colouring and type of interactions must be explained clearly to avoid ambiguity. {\color{red}{Additionally, the docking figures should be consistent, and the orientation of the binding poses must clearly illustrate the molecular interactions. This approach will also enable straightforward comparisons of docking images within a congeneric series. Whenever feasible, the docking results for theoretical compounds should be discussed alongside the known molecules, as this will enhance confidence in the findings.

Following are our recommendations while preparing the figures to illustrate the results in either 3-dimensional or 2-dimensional interactions. 
\begin{enumerate}
    \item How to Visualize 3D Interactions?:
    \begin{itemize}
	       \item Use molecular visualization tools like PyMOL\cite{bib65}, Chimera\cite{bib66}, or NGL Viewer\cite{bib67} to display protein-ligand complexes in 3D. 
           item Highlight key interactions such as hydrogen bonds, hydrophobic contacts, and $\pi$-stacking.
	       \item Label interacting residues and measure distances between ligand atoms and receptor residues to provide quantitative insights.
    \end{itemize}
    \item How to generate 2D Interaction Diagrams?:
        \begin{itemize}
            \item Tools like LigPlot+\cite{bib68}, PoseView\cite{bib69}, and Protein-Ligand Interaction Profiler (PLIP)\cite{bib70} can create clear 2D diagrams. 
            \item Display interactions such as hydrogen bonds and hydrophobic contacts with same orientation for intuitive comparison. 
        \end{itemize}    
\end{enumerate}    
}}

